\section{在格点上的实现}\label{sec:lattice}

要在格点上定义重正化的胶子准PDF，需要非微扰地确定重正化因子 $Z_{11}$、$Z_{22}$ 和 ${\overline {\delta m}}$。它们可以分别计算~\cite{Zhang:2017bzy}；也可以像夸克情形那样~\cite{Stewart:2017tvs,Chen:2017mzz}，在正规化无关的动量减除方案（RI/MOM）中把重正化因子作为整体算出~\cite{Martinelli:1994ty}。

在 RI/MOM 方案中，胶子准PDF的重正化因子可以这样确定：在一个远离壳的胶子态中计算相应算符的截腿格林函数，并要求抵消项在某个特定胶子动量处消去该格林函数中的所有圈图贡献。示意性地可以写成（RI/MOM 重正化的详细研究将另文给出~\cite{Todo:2018glu}）
\begin{align}
&\hspace*{-1.5em}[e^{{\overline{\delta m}}|z|}Z_{11} ^2 Z_3]^{-1}(z n,p^R_z, 1/a, \mu_R)=\non\\
&\hspace*{-2em}\left.\frac{P_{ij}^{ab} \langle 0| T [A^{a,i} (p)  O^{1}(z,0)A^{b,j}(-p) ]|0\rangle|_{amp}}{P_{ij}^{ab} \langle 0|T [A^{a,i} (p)   O^{1}(z,0)A^{b,j} (-p)]|0\rangle_{amp,tree}}\right|_{\tiny\begin{matrix}p^2=-\mu_R^2 \\ \!\!\!\!p_z=p^R_z\end{matrix}},
\end{align}
这是对 $O^1$ 的情形。这里 $A^{a,i}(p)$ 表示动量为 $p$ 的动量空间胶子场。$Z_3$ 是胶子场的重正化常数。对于 $O^{2,3,4}$，$Z_{11} ^2$ 应分别替换为 $Z_{22} ^2$、$Z_{11}Z_{12}$ 和 $Z_{22} ^2$。$P_{ij}^{ab}$ 是带色指标 $a,b$ 与洛伦兹指标 $i,j$ 的投影算符。投影的一个简单例子是 $P_{ij}^{ab}= \delta^{ab}g_{\perp,ij}$。
$\mu_R$ 和 $p_z^R$ 是 RI/MOM 方案中引入的非物理标度~\cite{Stewart:2017tvs,Chen:2017mzz}。

完成非微扰重正化之后，我们可以写出重正化胶子/夸克准PDF的因子化公式：
\begin{align}\label{allfac}
&\tilde g(x ,P_z, p_z^R, \mu_R)=\int_0^1\frac{dy}{y}\Big[C_{gg}\Big(\frac{x}{y}, \frac{\mu_R}{p_z^R},\frac{y P_z}{\mu}, \frac{y P_z}{p_z^R}\Big)g(y, \mu)\non\\
&+C_{gq}\Big(\frac{x}{y}, \frac{\mu_R}{p_z^R},\frac{y P_z}{\mu}, \frac{y P_z}{p_z^R}\Big)q_j(y, \mu)\Big]+\mathcal O\Big(\frac{M^2}{P_z^2}, \frac{\Lambda_{\rm QCD}^2}{P_z^2}\Big),\nonumber\\
&\tilde q_i(x ,P_z, p_z^R, \mu_R)=\int_0^1\frac{dy}{y}\Big[C_{q_i q_j}\Big(\frac{x}{y}, \frac{\mu_R}{p_z^R},\frac{y P_z}{\mu}, \frac{y P_z}{p_z^R}\Big)q_j(y, \mu)\non\\
&+C_{qg}\Big(\frac{x}{y}, \frac{\mu_R}{p_z^R},\frac{y P_z}{\mu}, \frac{y P_z}{p_z^R}\Big)g(y, \mu)\Big]+\mathcal O\Big(\frac{M^2,}{P_z^2}, \frac{\Lambda_{\rm QCD}^2}{P_z^2}\Big),
\end{align}
其中所有部分子分布的动量分数都处于 $[0,1]$ 区间，并且隐含对 $j$ 遍历所有夸克/反夸克味的求和。与重正化准PDF一样，硬系数 $C_{gg}$、$C_{qg}$、$C_{q_i q_j}$ 和 $C_{gq}$ 也依赖于非物理标度，例如 $\mu_R$、$p_z^R$ 以及强子动量 $P_z$。但从式~(\ref{allfac}) 提取出的 PDFs 不依赖这些标度。对于极化准PDF，因子化具有与式~(\ref{allfac}) 类似的形式。
利用算符乘积展开~\cite{Izubuchi:2018srq} 证明上述因子化公式的细节，以及硬系数单圈精度的微扰计算，将在另一篇文章中给出~\cite{Todo:2018glu}。
